\documentclass[10pt]{article}
\usepackage[T1]{fontenc}
\usepackage[utf8]{inputenc}
\usepackage[french]{babel}
\usepackage{amsmath,amssymb,mathrsfs,mathpazo}
\usepackage{geometry}
\geometry{a4paper,margin=2cm}
\usepackage{graphicx,array,tabularx,tabularray,colortbl,multirow,multicol,booktabs,enumitem,pifont,ulem,fancybox,diagbox,cancel}
\usepackage{pstricks,pstricks-add,pst-plot,pst-tree,pst-node,pst-text,pst-func,pst-math,pst-bspline,pst-3dplot,pst-eucl,pst-solides3d}
\let\pscalculate\undefined
\usepackage{pst-bezier}
\usepackage{tikz,pgfplots,tkz-tab}
\usetikzlibrary{arrows,arrows.meta,patterns,calc,positioning,trees,decorations.markings,shapes.geometric,shapes.misc,shapes.arrows,intersections,decorations.pathmorphing,decorations.pathreplacing,backgrounds,shadows,babel}
\pgfplotsset{compat=1.18}
\usepgfplotslibrary{fillbetween}
\usepackage[np]{numprint}
\usepackage{listings,makecell,pas-tableur,twoopt}
\let\barre\undefined
\usepackage[pstricks]{tablvar}
\definecolor{darkgreen}{RGB}{0,100,0}
\providecommand{\red}{\color{red}}
\providecommand{\blue}{\color{blue}}
\providecommand{\green}{\color{green}}
\providecommand{\black}{\color{black}}
\providecommand{\white}{\color{white}}
\providecommand{\gray}{\color{gray}}
\providecommand{\cyan}{\color{cyan}}
\providecommand{\magenta}{\color{magenta}}
\providecommand{\yellow}{\color{yellow}}
\providecommand{\violet}{\color{violet}}
\providecommand{\orange}{\color{orange}}
\providecommand{\euro}{\texteuro}
\providecommand{\EUR}{\texteuro}
\providecommand{\ieme}{\textsuperscript{e}}
\providecommand{\degres}{\ensuremath{^\circ}}
\providecommand{\degree}{\ensuremath{^\circ}}
\providecommand{\vv}[1]{\overrightarrow{#1}}
\providecommand{\Coord}[1]{#1}
\providecommand{\decosix}{}
\providecommand{\starredbullet}{\ensuremath{\bullet}}
\pagestyle{empty}
\setlength{\parindent}{0pt}
\usepackage{adjustbox}\DeclareUnicodeCharacter{2236}{:}\DeclareUnicodeCharacter{2190}{\ensuremath{\leftarrow}}
\providecommand{\R}{\mathbb{R}}
\providecommand{\N}{\mathbb{N}}
\providecommand{\Z}{\mathbb{Z}}
\providecommand{\Q}{\mathbb{Q}}
\providecommand{\D}{\mathbb{D}}
\providecommand{\C}{\mathbb{C}}
\providecommand{\e}{{\mathrm{e}}}
\providecommand{\vect}[1]{\overrightarrow{#1}}
\providecommand{\vectt}[1]{\overrightarrow{\mathrm{#1}}}
\providecommand{\vvt}[1]{\overrightarrow{\mathrm{#1}}}
\providecommand{\np}[2][]{{#2~#1}}
\providecommand{\vv}[1]{{\overrightarrow{#1}}}
\renewcommand{\d}{{\mathrm{d}}}
\providecommand{\degres}{{^\circ}}
\providecommand{\degre}{{^\circ}}
\providecommand{\blue}{\color{blue}}
\providecommand{\red}{\color{red}}
\providecommand{\enskip}{\quad}
\providecommand{\llbracket}{\lbrack\!\lbrack}
\providecommand{\rrbracket}{\rbrack\!\rbrack}
\providecommand{\euro}{{\text{€}}}
\providecommand{\no}{{\text{n°}}}
\providecommand{\interval}[2]{{[#1;#2]}}
\providecommand{\intervFO}[2]{{[#1;#2[}}
\providecommand{\dotfill}{\ldots}

\providecommand{\ds}{\displaystyle}
\providecommand{\textding}[1]{#1}
\providecommand{\Oij}{\(\mathrm{O};\vec{\imath},\vec{\jmath}\)}
\providecommand{\Oijk}{\(\mathrm{O};\vec{\imath},\vec{\jmath},\vec{k}\)}
\providecommand{\Ouv}{\(\mathrm{O};\vec{u},\vec{v}\)}
\providecommand{\intervFF}[2]{\ensuremath{[#1;#2]}}
\renewcommand{\euro}{\ifmmode\text{\texteuro}\else\texteuro\fi}
\DeclareUnicodeCharacter{00B0}{\ifmmode^{\circ}\else\textdegree\fi}
\DeclareUnicodeCharacter{20AC}{\ifmmode\text{\texteuro}\else\texteuro\fi}
% Énoncés : Éducation nationale. Transcription et corrigés : APMEP, Denis Vergès et contributeurs. Découpage et classement : CoopMaths, Sébastien Lozano, Éric Elter et contributeurs. Adaptation de la présentation : Maths Explorer.
% CC BY-SA : https://coopmaths.fr/www/about/
% https://coopmaths.fr/alea/static/bac/2025/tex/bac_2025_11_sujet2_ameriquesud_3_cor.tex
\begin{document}
\noindent\textit{APMEP / CoopMaths. CC BY-SA. Présentation adaptée par Maths Explorer.}\par\medskip
\medskip


On considère les suites $(v_n)$ et $(w_n)$ définies pour tout entier naturel $n$
par :
\[ \left\{\begin{array}{l c l}
v_0 &=& \ln (4)\\
v_{n+1} &=& \ln \left(- 1 +2\e^{v_n}\right)
\end{array}\right. \qquad \text{et} \qquad w_n = - 1 + \e^{v_n}.\]

On admet que la suite $(v_n)$ est bien définie et strictement positive.












































































\begin{enumerate}
\item $\bullet~~$$v_1 = \ln \left(- 1 + 2\e^{v_0}\right) = \ln \left(- 1 + 2\e^{\ln 4}\right) = \ln (- 1 + 2 \times 4) = \ln (- 1 + 8)  = \ln 7$ (valeur approchée ligne 4 colonne 3)

$\bullet~~$$w_0 = - 1 + \e^{v_0} = - 1 + \e^{\ln 4} = - 1 + 4 = 3$ (valeur ligne 3 colonne 4)
\item
	\begin{enumerate}
		\item Il faut saisir la formule 2.
		\item On peut penser que la suite $(v_n)$ est croissante
		\item Démonstration par récurrence :

\emph{Initialisation} : on a $v_0 = \ln 4$ et $v_1 = \ln 7$ : la croissance de la fonction ln assure que $v_0 < v_1$.

\emph{Hérédité} : on suppose que pour $n \in \N, \:v_n < v_{n+1}$, alors :

$\e^{v_n} < \e^{v_{n+1}}$ par croissance de la fonction exponentielle

$2\e^{v_n} < 2\e^{v_{n+1}}$ la multiplication par $+ 2$ respecte l'ordre

$- 1 + 2\e^{v_n} < - 1 + 2\e^{v_{n+1}}$ l'addition respecte l'ordre

$\ln \left(- 1 + 2\e^{v_n}\right) < \ln \left(- 1 + 2\e^{v_{n+1}}\right)$ par croissance de la fonction logarithme népérien (on a supposé que la suite $(v_n)$ est bien définie donc que tous les nombres de la forme $- 1 + 2\e^{v_n}$ sont supérieurs à zéro) ; soit finalement 

$v_{n+1} < v_{n+2}$ : l'inégalité est vraie au rang $n + 1$.

Conclusion : l'inégalité est vraie au rang 0 et si elle est vraie au rang $n \in \N$ elle l'est aussi au rang suivant : d'après le principe de récurrence :
\begin{center} quel que soit $n \in \N, \quad v_n < v_{n+1}$ \end{center}

La suite $(v_n)$ est croissante.
	\end{enumerate}
\item
	\begin{enumerate}
		\item Quel que soit $n \in \N$, \: $w_n = - 1 + \e^{v_n}$, donc 
		
$w_{n+1} = - 1 + \e^{v_{n+1}} = - 1 + \e^{\ln \left(- 1 + 2\e^{v_n}\right)} = - 1 - 1 + 2\e^{v_n} = -2 + 2\e^{v_n} = 2\left(- 1 + \e^{v_n}\right)$ : finalement 

Quel que soit $n \in \N$, \quad $w_{n+1} = 2w_n$ : cette égalité montre que la suite $w_n$ est géométrique de raison 2 avec pour premier terme $w_0 = 3$.
		\item On sait qu'alors le terme général $w_n$ est égal à :
		
		\begin{center} $w_n = 3 \times 2^n$ \end{center}
		
Or par définition $w_n = - 1 + \e^{v_n} = 3 \times 2^n \iff \e^{v_n} = 1 + 3 \times 2^n$.

Par croissance de la fonction logarithme népérien : $v_n = \ln \left( 1 + 3 \times 2^n\right)$ pour tout entier naturel.
\item On sait que $\displaystyle\lim_{n \to + \infty} 2^n = + \infty$, donc $\displaystyle\lim_{n \to + \infty} 3 \times 2^n = + \infty$, puis $\displaystyle\lim_{n \to + \infty} 1 + 3 \times 2^n = + \infty$ et enfin $\displaystyle\lim_{n \to + \infty} \ln \left(1 + 3 \times 2^n\right)= + \infty$.

Conclusion : $\displaystyle\lim_{n \to + \infty} v_n = + \infty$.
	\end{enumerate}
\item L'algorithme permet de calculer les termes de la suite $v_n$ à partir de $v_0$. On vient de démontrer que la suite $(v_n)$ n'est pas majorée. Donc quel que soit le choix du nombre $S$, il existe un rang $p$, tel que $v_p > S$ et l'algorithme donnera ce rang $p$.
\end{enumerate}

\medskip

\end{document}
